\documentclass[a4paper,12pt]{article}
\usepackage{amsmath,amssymb,amsthm,hyperref}
\ProvidesPackage{journal}
\usepackage{hyperref}
\usepackage[mathcal,mathbf]{euler}
\usepackage{enumitem}
\usepackage{booktabs}
\usepackage{adjustbox}
\usepackage{mathtools}
\usepackage{listings}
\usepackage{amsfonts}
\usepackage{amscd}
\usepackage{amsmath}
\usepackage{amssymb}
\usepackage{amsthm}
\usepackage{tikz}
\usepackage{tikz-cd}
\usepackage{url}
\usepackage[utf8]{inputenc}
\usepackage[english,ukrainian]{babel}
%\usepackage[T1]{fontenc}
\usepackage[only,llbracket,rrbracket,llparenthesis,rrparenthesis]{stmaryrd}

\usetikzlibrary{babel}

\newcommand*{\incmap}{\hookrightarrow}
\newcommand*{\thead}[1]{\multicolumn{1}{c}{\bfseries #1}}
\renewcommand{\Join}{\vee} % Join operation symbol
\newcommand{\tabstyle}[0]{\scriptsize\ttfamily\fontseries{l}\selectfont}

\newcommand{\Mod}{\text{Mod}}
\newcommand{\dg}{\text{dg}}
\newcommand{\Z}{\mathbb{Z}}
\newcommand{\Q}{\mathbb{Q}}
\newcommand{\R}{\mathbb{R}}
\newcommand{\RP}{\mathbb{RP}}
\newcommand{\CP}{\mathbb{CP}}
\newcommand{\Tor}{\text{Tor}}
\newcommand{\Ext}{\text{Ext}}
\newcommand{\Hom}{\text{Hom}}
\newcommand{\id}{\text{id}}
\newcommand{\pt}{\text{pt}}
\newcommand{\im}{\text{im}}
\newcommand{\Ab}{\text{Ab}}
\newcommand{\Ch}{\text{Ch}}
\newcommand{\Spectra}{\text{Spectra}}
\newcommand{\Ssphere}{\mathbb{S}}
\newcommand{\HA}{H}
\newcommand{\KU}{\text{KU}}
\newcommand{\KO}{\text{KO}}
\newcommand{\T}{\mathbb{T}}
\newcommand{\Aone}{\mathbb{A}^1}
\newcommand{\Sm}{\text{Sm}}
\newcommand{\Spt}{\text{Sp}}
\newcommand{\Nis}{\text{Nis}}

\lstset{basicstyle=\footnotesize,inputencoding=utf8,extendedchars=true,literate={
{𝟎}{{\ensuremath{\mathbf{0}}}}1
{𝟏}{{\ensuremath{\mathbf{1}}}}1
{𝟐}{{\ensuremath{\mathbf{2}}}}1
{≔}{{\ensuremath{\mathrm{:=}}}}1
{α}{{\ensuremath{\mathrm{\alpha}}}}1
{ᵂ}{{\ensuremath{^W}}}1
{β}{{\ensuremath{\mathrm{\beta}}}}1
{γ}{{\ensuremath{\mathrm{\gamma}}}}1
{δ}{{\ensuremath{\mathrm{\delta}}}}1
{ε}{{\ensuremath{\mathrm{\varepsilon}}}}1
{ζ}{{\ensuremath{\mathrm{\zeta}}}}1
{η}{{\ensuremath{\mathrm{\eta}}}}1
{θ}{{\ensuremath{\mathrm{\theta}}}}1
{ι}{{\ensuremath{\mathrm{\iota}}}}1
{κ}{{\ensuremath{\mathrm{\kappa}}}}1
{λ}{{\ensuremath{\mathrm{\lambda}}}}1
{μ}{{\ensuremath{\mathrm{\mu}}}}1
{ν}{{\ensuremath{\mathrm{\nu}}}}1
{ξ}{{\ensuremath{\mathrm{\xi}}}}1
{π}{{\ensuremath{\mathrm{\mathnormal{\pi}}}}}1
{ρ}{{\ensuremath{\mathrm{\rho}}}}1
{σ}{{\ensuremath{\mathrm{\sigma}}}}1
{τ}{{\ensuremath{\mathrm{\tau}}}}1
{φ}{{\ensuremath{\mathrm{\varphi}}}}1
{χ}{{\ensuremath{\mathrm{\chi}}}}1
{ψ}{{\ensuremath{\mathrm{\psi}}}}1
{ω}{{\ensuremath{\mathrm{\omega}}}}1
{Π}{{\ensuremath{\mathrm{\Pi}}}}1
{П}{{\ensuremath{\mathrm{\Pi}}}}1
{Γ}{{\ensuremath{\mathrm{\Gamma}}}}1
{Δ}{{\ensuremath{\mathrm{\Delta}}}}1
{Θ}{{\ensuremath{\mathrm{\Theta}}}}1
{Λ}{{\ensuremath{\mathrm{\Lambda}}}}1
{Σ}{{\ensuremath{\mathrm{\Sigma}}}}1
{Φ}{{\ensuremath{\mathrm{\Phi}}}}1
{Ξ}{{\ensuremath{\mathrm{\Xi}}}}1
{Ψ}{{\ensuremath{\mathrm{\Psi}}}}1
{Ω}{{\ensuremath{\mathrm{\Omega}}}}1
{ℵ}{{\ensuremath{\aleph}}}1
{≤}{{\ensuremath{\leq}}}1
{≥}{{\ensuremath{\geq}}}1
{≠}{{\ensuremath{\neq}}}1
{≈}{{\ensuremath{\approx}}}1
{≡}{{\ensuremath{\equiv}}}1
{≃}{{\ensuremath{\simeq}}}1
{≤}{{\ensuremath{\leq}}}1
{≥}{{\ensuremath{\geq}}}1
{∂}{{\ensuremath{\partial}}}1
{∆}{{\ensuremath{\triangle}}}1 % or \laplace?
{∫}{{\ensuremath{\int}}}1
{∑}{{\ensuremath{\mathrm{\Sigma}}}}1
{→}{{\ensuremath{\rightarrow}}}1
{⊥}{{\ensuremath{\perp}}}1
{∞}{{\ensuremath{\infty}}}1
{∂}{{\ensuremath{\partial}}}1
{∓}{{\ensuremath{\mp}}}1
{±}{{\ensuremath{\pm}}}1
{×}{{\ensuremath{\times}}}1
{⊕}{{\ensuremath{\oplus}}}1
{⊗}{{\ensuremath{\otimes}}}1
{⊞}{{\ensuremath{\boxplus}}}1
{∇}{{\ensuremath{\nabla}}}1
{√}{{\ensuremath{\sqrt}}}1
{⬝}{{\ensuremath{\cdot}}}1
{•}{{\ensuremath{\cdot}}}1
{∘}{{\ensuremath{\circ}}}1
{⁻}{{\ensuremath{^{-}}}}1
{▸}{{\ensuremath{\blacktriangleright}}}1
{★}{{\ensuremath{\star}}}1
{∧}{{\ensuremath{\wedge}}}1
{∨}{{\ensuremath{\vee}}}1
{¬}{{\ensuremath{\neg}}}1
{⊢}{{\ensuremath{\vdash}}}1
{⟨}{{\ensuremath{\langle}}}1
{⟩}{{\ensuremath{\rangle}}}1
{↦}{{\ensuremath{\mapsto}}}1
{→}{{\ensuremath{\rightarrow}}}1
{↔}{{\ensuremath{\leftrightarrow}}}1
{⇒}{{\ensuremath{\Rightarrow}}}1
{⟹}{{\ensuremath{\Longrightarrow}}}1
{⇐}{{\ensuremath{\Leftarrow}}}1
{⟸}{{\ensuremath{\Longleftarrow}}}1
{∩}{{\ensuremath{\cap}}}1
{∪}{{\ensuremath{\cup}}}1
{·}{{\ensuremath{\cdot}}}1
{ᵢ}{{\ensuremath{_i}}}1
{ⱼ}{{\ensuremath{_j}}}1
{₊}{{\ensuremath{_+}}}1
{ℑ}{{\ensuremath{\Im}}}1
{𝒢}{{\ensuremath{\mathcal{G}}}}1
{ℕ}{{\ensuremath{\mathbb{N}}}}1
{𝟘}{{\ensuremath{\mathbb{0}}}}1
{ℤ}{{\ensuremath{\mathbb{Z}}}}1
{ℝ}{{\ensuremath{\mathbb{R}}}}1
{ℙ}{{\ensuremath{\mathbb{P}}}}1
{∁}{{\ensuremath{\subseteq}}}1
{⊂}{{\ensuremath{\subseteq}}}1
{⊆}{{\ensuremath{\subseteq}}}1
{⊄}{{\ensuremath{\nsubseteq}}}1
{⊈}{{\ensuremath{\nsubseteq}}}1
{⊃}{{\ensuremath{\supseteq}}}1
{⊇}{{\ensuremath{\supseteq}}}1
{⊅}{{\ensuremath{\nsupseteq}}}1
{⊉}{{\ensuremath{\nsupseteq}}}1
{∈}{{\ensuremath{\in}}}1
{∉}{{\ensuremath{\notin}}}1
{∋}{{\ensuremath{\ni}}}1
{∌}{{\ensuremath{\notni}}}1
{∅}{{\ensuremath{\emptyset}}}1
{∖}{{\ensuremath{\setminus}}}1
{†}{{\ensuremath{\dag}}}1
{ℕ}{{\ensuremath{\mathbb{N}}}}1
{ℤ}{{\ensuremath{\mathbb{Z}}}}1
{ℝ}{{\ensuremath{\mathbb{R}}}}1
{ℚ}{{\ensuremath{\mathbb{Q}}}}1
{ℂ}{{\ensuremath{\mathbb{C}}}}1
{⌞}{{\ensuremath{\llcorner}}}1
{⌟}{{\ensuremath{\lrcorner}}}1
{⦃}{{\ensuremath{ \{\!| }}}1
{⦄}{{\ensuremath{ |\!\} }}}1
{ᵁ}{{\ensuremath{^U}}}1
{₋}{{\ensuremath{_{-}}}}1
{₁}{{\ensuremath{_1}}}1
{₂}{{\ensuremath{_2}}}1
{₃}{{\ensuremath{_3}}}1
{₄}{{\ensuremath{_4}}}1
{₅}{{\ensuremath{_5}}}1
{₆}{{\ensuremath{_6}}}1
{₇}{{\ensuremath{_7}}}1
{₈}{{\ensuremath{_8}}}1
{₉}{{\ensuremath{_9}}}1
{₀}{{\ensuremath{_0}}}1
{¹}{{\ensuremath{^1}}}1
{ₙ}{{\ensuremath{_n}}}1
{⊤}{{\ensuremath{\top}}}1
{↪}{{\ensuremath{\hookrightarrow}}}1
{⇒}{{\ensuremath{\Rightarrow}}}1
{ₘ}{{\ensuremath{_m}}}1
{↑}{{\ensuremath{\uparrow}}}1
{↓}{{\ensuremath{\downarrow}}}1
{▸}{{\ensuremath{\triangleright}}}1
{∀}{{\ensuremath{\forall}}}1
{∃}{{\ensuremath{\exists}}}1
{λ}{{\ensuremath{\mathrm{\lambda}}}}1
{=}{{\ensuremath{=}}}1
{<}{{\ensuremath{\textless}}}1
{>}{{\ensuremath{\textgreater}}}1
{_}{{$\_$}}1
{‖}{{\ensuremath{\Vert}}}1
{→}{{$\rightarrow$}}1
{Π}{{$\Pi$}}1
{Σ}{{$\Sigma$}}1
{Г}{{$\Gamma$}}1
{Θ}{{$\Theta$}}1
{∆}{{$\Delta$}}1
{Σ}{{$\Sigma$}}1
{Ф}{{$\Phi$}}1
{σ}{{$\sigma$}}1
{δ}{{$\delta$}}1
{ν}{{$\nu$}}1
{η}{{$\eta$}}1
{₁}{{$_1$}}1
{•}{{$\bullet$}}1
{є}{{$\varepsilon$}}1
{ᵀ}{{$\textsuperscript{T}$}}1
{ᵗ}{{$\textsuperscript{t}$}}1
{◊}{{$\lozenge$}}1
{ᵀ}{{$^T$}}1
{ᵗ}{{$^t$}}1
 }
}

\theoremstyle{definition}

\newtheorem{definition}{Definition}
\newtheorem{theorem}{Theorem}
\newtheorem{lemma}{Lemma}
\newtheorem{notation}{Notation}
\newtheorem{acknowledgment}{Acknowledgment}
\newtheorem{example}{Example}
\newtheorem{remark}{Remark}
\newtheorem{corollary}{Corollary}
\newtheorem{proposition}{Proposition}

\newcommand{\Spec}{\mathbf{Spec}}

\def\mapright#1{\xrightarrow{{#1}}}
\def\mapleft#1{\xleftarrow{{#1}}}
\def\under{\underline{\hspace{0.25cm}}}

\lefthyphenmin=1
\hyphenpenalty=100
\tolerance=11000
\emergencystretch=1em
\hfuzz=2pt
\vfuzz=2pt

\newif\ifincludeTOC
\includeTOCtrue

\def\mapright#1{\xrightarrow{{#1}}}
\def\mapleft#1{\xleftarrow{{#1}}}
\def\mapup#1{\Big\uparrow\rlap{\raise2pt{\scriptstyle{#1}}}}
\def\mapupl#1{\Big\uparrow\llap{\raise2pt{\scriptstyle{#1}}}}
\def\mapdown#1{\Big\downarrow\rlap{\raise2pt{\scriptstyle{#1}}}}
\def\mapdownl#1{\Big\downarrow\llap{\raise2pt{\scriptstyle{#1}}}}
\def\mapdiagl#1{\vcenter{\searrow}\rlap{\raise2pt{\scriptstyle{#1}}}}
\def\mapdiagr#1{\vcenter{\swarrow}\rlap{\raise2pt{\scriptstyle{#1}}}}
\def\under{\underline{\hspace{0.25cm}}}


\begin{document}

\title{Issue XXXVII: The Ladder of States}
\author{Namdak Tonpa}
\date{\today}
\maketitle

\begin{abstract}
У цій роботі ми розглядаємо «драбину станів» — ієрархію математичних структур, що описують стани та їх еволюцію, від гомотопічної теорії типів через класичну ймовірність і марковські процеси до квантової механіки. 
Кожен щабель інтерпретується як перехід між категоріями, де стани є узагальненими елементами, а динаміка — морфізмами. 
Особлива увага приділяється категоріальній лінеаризації марковської динаміки в гільбертових просторах і аналітичному зв’язку між стохастичними та квантовими еволюціями через Wick rotation та формулу Фейнмана–Каца.
\end{abstract}

\ifincludeTOC
\tableofcontents
\fi

\section{Категорія станів}
\subsection*{Стан як об’єкт, еволюція як морфізм}

У рамках підходу \emph{Groupoid Infinity} математика розглядається як єдина суверенна система, в якій фізичні та ймовірнісні структури моделюються як об’єкти у відповідних $\infty$-категоріях, а динамічні процеси — як морфізми (1-морфізми у вищому сенсі).

«Драбина станів» від $L^2$-просторів до квантової механіки не є простою послідовністю прикладів. Це \textbf{функторіальна драбина} — ланцюг функторів між категоріями, де кожен щабель відповідає переходу між різними типами просторів станів та їх еволюціями.

Ми формалізуємо базовий принцип:
\begin{itemize}
    \item \textbf{Стан} — це узагальнений елемент простору (морфізм $1 \to X$);
    \item \textbf{Еволюція} — це оператор (морфізм або ендофунктор).
\end{itemize}

Різні фізичні теорії виникають як різні класи просторів і дозволених морфізмів. У HoTT/$L^2$ стає об’єктом у $\dagger$-категорії $\mathbf{Hilb}$, марковські процеси породжують монаду ймовірності на категорії вимірних просторів, а самі процеси реалізуються як морфізми у відповідній Kleisli-категорії. Квантова механіка постає як синтез у dagger-структурі.

Нижче ми крок за кроком піднімаємося драбиною, а в кінці наводимо категоріальну діаграму переходів.

\subsection{Мова нескінченновимірних станів}

Функціональний аналіз є базовою мовою для роботи з Banach- і Hilbert-просторами. Центральним об’єктом є гільбертів простір $L^2(X,\mu)$ квадратно-інтегровних функцій з внутрішнім добутком
\[
\langle f \mid g \rangle = \int_X f(x) \overline{g(x)}\,d\mu(x).
\]
Кожен елемент $\psi \in L^2(X,\mu)$ уже інтерпретується як \textbf{вектор стану}. Обмежені оператори є ендоморфізмами, а часова еволюція описується сильною неперервною напівгрупою операторів на $L^p$-просторах (зокрема на $L^2$ за умов інваріантності міри).

$L^2$-простори природно утворюють $\dagger$-категорію $\mathbf{Hilb}$ (категорію гільбертових просторів з обмеженими операторами та інволюцією $\dagger$, що відповідає спряженому оператору).

\subsection{Rigged Hilbert space як узагальнення}

Звичайні $L^2$-функції недостатні для опису $\delta$-функцій Дірака, власних функцій неперервного спектру чи сингулярних розв’язків диференціальних рівнянь. Теорія розподілів Шварца вводить \textbf{rigged Hilbert space} (трійку Гельфанда):
\[
\mathcal{S} \subset L^2 \subset \mathcal{S}',
\]
де $\mathcal{S}$ — ядерний простір тест-функцій (простір Шварца), а $\mathcal{S}'$ — простір розподілів (його топологічний дуал).

У категоріальній мові розподіли можна розглядати як узагальнені елементи подвійного об’єкта, що узгоджується з ідеєю generalized elements у сенсі Гротендіка. Саме в $\mathcal{S}'$ «живуть» плоскі хвилі та інші сингулярні об’єкти, необхідні для спектральної теорії. 

Rigged Hilbert space є першим кроком до \textbf{cohesive modalities} у синтетичній диференціальній геометрії: розподіли поводяться як «інфінітезимальні» або «проінфінітні» елементи.

\subsection{$L^2$ як простір ймовірностей}

У теорії ймовірностей випадкова величина $X \in L^2(\Omega, \mathcal{F}, \mathbb{P})$ задовольняє
\[
\|X\|_2^2 = \mathbb{E}[X^2].
\]
Ймовірнісний простір $(\Omega, \mathcal{F}, \mathbb{P})$ індукує гільбертів простір $L^2(\Omega, \mathbb{P})$, в якому умовне сподівання $\mathbb{E}[\,\cdot\,\mid\mathcal{G}]$ є ортогональною проєкцією на підпростір $L^2(\Omega, \mathcal{G}, \mathbb{P})$.

Незалежність $\sigma$-алгебр імплікує ортогональність відповідних центрованих підпросторів у $L^2$, хоча й не зводиться до неї повністю. Таким чином, значна частина класичної теорії ймовірностей має природну інтерпретацію в термінах функціонального аналізу.

Категоріально ймовірнісна міра відповідає стану (positive normalized functional) на комутативній $C^*$-алгебрі $L^\infty(\Omega)$.

\subsection{Категорії Маркова}

Ймовірнісна динаміка природно організовується на рівні категорії процесів. Нехай $\mathbf{Meas}$ — категорія вимірних просторів. Монада Гірі ймовірнісних розподілів
\[
\mathcal{G} : \mathbf{Meas} \to \mathbf{Meas}
\]
індукує Kleisli-категорію $\mathbf{Kl}(\mathcal{G})$, яку ми позначаємо $\mathbf{Markov}$. 

У категорії $\mathbf{Markov}$:
\begin{itemize}
    \item об’єкти — вимірні простори;
    \item морфізми $X \to Y$ — марковські ядра $X \to \mathcal{G}(Y)$;
    \item стани — морфізми $1 \to X$;
    \item еволюції — морфізми $X \to Y$.
\end{itemize}

Умовна ймовірність і правило Байєса набувають природної категоріальної інтерпретації через універсальні властивості (заповнення діаграм).

\subsection{Марковські процеси як напівгрупи операторів}

Марковський процес можна розглядати як одно-параметричну напівгрупу ендоморфізмів у $\mathbf{Markov}$: $T_t : X \to X$. 

У функціональному аналізі цьому відповідає напівгрупа операторів на $L^2(X,\mu)$:
\[
(T_t f)(x) = \mathbb{E}[f(X_t) \mid X_0 = x].
\]
Її інфінітезимальний генератор $L$ задовольняє рівняння Колмогорова–Чепмена. 

Таким чином, марковська динаміка реалізується як контрактна додатно збереженна напівгрупа на гільбертовому просторі — структура, дуже близька до квантової механіки, але з контрактними замість унітарних операторів.

Простір траєкторій марковського процесу природно моделюється як об’єкт у $\infty$-категорії типів. Функтор $\Pi_\infty : \mathbf{Top} \to \infty\text{-}\mathbf{Gpd}$ піднімає топологічну модель до гомотопічної, де шляхи відповідають вищим гомотопіям.

\subsection{Геометризація фазового простору}

У класичній механіці фазовий простір $M = T^*Q$ оснащений симплектичною формою $\omega$. Гамільтонів потік індукує еволюцію на алгебрі гладких функцій $C^\infty(M)$ через оператор Лі. 

Перехід до $L^2(M, \mu)$ (з мірою Ліувіля) та подальше розширення до rigged Hilbert space є ключовим кроком до \textbf{geometric quantization} — природного мосту між класичною та квантовою механікою.

\subsection{Нелінійні системи}

Для нелінійних систем $\dot{x} = F(x)$ еволюція стає нелінійним потоком. У функціональному аналізі це відповідає нелінійним напівгрупам, теорії хаосу та біфуркацій. 

У HoTT такі траєкторії природно моделюються за допомогою higher inductive types або cohesive modalities, які дозволяють розглядати простір шляхів як $\infty$-групоїд.

\subsection{Квантова механіка}

У квантовій механіці стан описується вектором $\psi \in \mathcal{H}$, спостережувані — самоспряженими операторами $A = A^\dagger$, а еволюція задається унітарною групою:
\[
i\hbar \frac{d\psi}{dt} = H\psi \quad \Longrightarrow \quad U_t = e^{-iHt/\hbar}.
\]

Квантова теорія синтезує попередні рівні:
\begin{itemize}
    \item $L^2$ + rigged структура (спектральна теорія);
    \item ймовірність ($|\psi|^2$ як густина ймовірності);
    \item операторну динаміку (перехід від марковських до унітарних еволюцій);
    \item геометрію (симплектична геометрія $\to$ преквантовація).
\end{itemize}

\subsection{Обертання Віка та формула Фейнмана–Каца}

Глибокий аналітичний зв’язок між марковською та квантовою динамікою забезпечується Wick rotation ($t \mapsto -it$) та формулою Фейнмана–Каца. 

Формула Фейнмана–Каца виражає розв’язок рівняння Шредінгера через умовне сподівання відносно евклідового марковського процесу. Формально це відповідає аналітичному продовженню генератора дифузії $L$ до самоспряженого гамільтоніана $H$ (у багатьох моделях $H \approx -L$).

У $\dagger$-категоріях цей перехід інтерпретується як перехід від позитивних контрактних напівгруп до унітарних груп, хоча канонічного функторіального підняття загалом не існує.

\subsection{Категоріальна діаграма переходів}

У рамках HoTT/$\infty$-категорій драбина станів може бути представлена такою (спрощеною) комутативною діаграмою:

\[
\begin{CD}
\infty\text{-}\mathbf{Gpd} @>>> \mathbf{Top} @>>> \mathbf{Meas} \\
@. @. @VV{\mathcal{G}\ (\text{Giry})}V \\
@. @. \mathbf{Markov} @>{\mathcal{K}}>> \mathbf{Hilb} \\
@. @. @VV{\text{generator }L}V @VV{\text{rigging}}V \\
@. @. \mathbf{Semigroup} @. \mathbf{Quant} \\
@. @. @VV{\text{Wick / Feynman--Kac}}V \\
@. @. @. \mathbf{UnitaryGroup}
\end{CD}
\]

Додатково:
\[
\begin{CD}
\mathbf{Prob} @>{\text{cond.\ exp.}}>> \mathbf{Hilb} \\
@V{\text{Kolmogorov}}VV @VV{\text{rigging}}V \\
\mathbf{Markov} @>{\text{Feynman--Kac}}>> \mathbf{Quant} \\
@V{\text{nonlinear}}VV @VV{\text{geometric quantization}}V \\
\mathbf{NonlinearDyn} @. \mathbf{CohesiveModality}
\end{CD}
\]

\subsection{Функторіальна лінеаризація}

Ми визначаємо композицію функторів, яка реалізує лінеаризацію гомотопічної динаміки:
\[
\mathcal{L} = \mathcal{K} \circ \mathcal{G} \circ \mathcal{B} \circ |{-}| \colon \infty\text{-}\mathbf{Gpd} \to \mathbf{Hilb},
\]
де $\mathcal{K}$ — гільбертова реалізація марковських ядер.

На рівні марковської категорії функтор $\mathcal{K} : \mathbf{Markov} \to \mathbf{Hilb}$ діє так:
- на об’єктах: $X \mapsto L^2(X, \mu)$;
- на морфізмах (марковське ядро $K : X \to \mathcal{G}(Y)$):
\[
(\mathcal{K}K f)(x) = \int_Y f(y) \, K(x, dy).
\]

Отриманий оператор є лінійним, позитивним і контрактним.

\subsection*{Висновок}

Ми побудували функторіальну факторизацію динамічних систем:
\[
\infty\text{-}\mathbf{Gpd} \;\to\; \mathbf{Top} \;\to\; \mathbf{Meas} \;\to\; \mathbf{Markov} \;\to\; \mathbf{Hilb}.
\]

Ця драбина показує, що:
\begin{itemize}
    \item гомотопічні структури задають простори станів і траєкторій;
    \item вимірна структура вводить ймовірність;
    \item категорії Маркова описують стохастичну динаміку;
    \item гільбертові простори забезпечують її лінеаризацію та операторну реалізацію.
\end{itemize}

Функціональний аналіз і квантова механіка постають як природні представлення нижчих рівнів цієї ієрархії. Cohesive homotopy type theory пропонує єдиний синтетичний фреймворк, у якому всі щаблі драбини можуть бути описані одночасно.

\begin{thebibliography}{9}
\bibitem{VolI} Homotopy Type Theory: Univalent Foundations of Mathematics.
\bibitem{VolIV} Cohesive Topoi and Geometric Quantization.
\bibitem{Shulman} M. Shulman, Modal Homotopy Type Theory.
\bibitem{Urs} Higher Category Theory and Physics (various works).
\bibitem{FeynmanKac} R. P. Feynman, M. Kac, Interaction of Markov processes and quantum theory.
\end{thebibliography}

\end{document}